\documentclass[../../main.tex]{subfiles}
\begin{document}
% SOURCE p=221; printed=206; object=ch06-maximal-lp-introduction; region=middle
最后给出 Marcinkiewicz 定理的两个应用。第一个涉及\cref{definition:ch03-hardy-littlewood-maximal}中的 Hardy–Littlewood 极大算子
\begin{align*}
Hf(x)=\sup_{r>0}\frac1{m(B(r,x))}\int_{B(r,x)}|f(y)|\,\mathrm dy,
\qquad f\in L^1_{\mathrm{loc}}(\mathbb R^n).
\end{align*}
$H$ 显然次线性，并对所有 $f\in L^\infty$ 满足 $\|Hf\|_\infty\leqslant\|f\|_\infty$。而\cref{theorem:ch03-maximal}恰说明它是弱 $(1,1)$ 型的。因此得到：
% SOURCE p=221; printed=206; object=ch06-maximal-lp; number=6.35; region=middle
\begin{corollary}\label{corollary:ch06-maximal-lp}
存在常数 $C>0$，使得若 $1<p<\infty$、$f\in L^p(\mathbb R^n)$，则
\begin{align*}
\|Hf\|_p\leqslant C\frac p{p-1}\|f\|_p.
\end{align*}
\end{corollary}
% SUPPLEMENT: id=v14-ch06-maximal-lp-proof; parent=ch06-maximal-lp; kind=omitted-proof+motivation+hypotheses; origin=independent-derivation
\begin{proof}
\textbf{常数估计的补全（重构新增）。}记极大定理的弱 $(1,1)$ 型常数为 $A_n$。局部积分的三角不等式及上确界给出 $H(f+g)\leqslant Hf+Hg$，且 $H(cf)=cHf$ 对 $c>0$ 成立。$L^1+L^\infty\subset L^1_{\mathrm{loc}}$，故这些输入均在定义域中；其可测性已由\cref{proposition:ch03-maximal-measurable}保证。对本质有界 $f$，每个球上的平均不超过 $\|f\|_\infty$，所以强 $(\infty,\infty)$ 型常数可取一。应用前面已完整计算的第二个特例，
\begin{align*}
\|Hf\|_p\leqslant2\left(\frac{A_n}{p-1}\right)^{1/p}\|f\|_p.
\end{align*}
又 $A_n^{1/p}\leqslant\max(1,A_n)$，并且 $(p-1)^{1-1/p}\leqslant p$：当 $1<p\leqslant2$ 时左边不超过一，当 $p\geqslant2$ 时左边不超过 $p-1$。因此取 $C=2\max(1,A_n)$ 即得所述界。
\end{proof}
\begin{remark}
\textbf{动机与逐条件说明（重构新增）。}极大算子一般不能把 $L^1$ 有界地映回 $L^1$，但它的弱 $(1,1)$ 界与平凡的无穷端点界足以恢复所有严格中间指数的强估计。$p>1$ 来自低水平积分的收敛，并解释常数接近一时可能发散；\cref{exercise:ch06-maximal-indicator}将给出这种发散确实出现的例子。$p<\infty$ 是本推论给出的范围；在 $p=\infty$ 已有常数一的直接估计，不应把它当作从公式代入无穷得到。空间为 $\mathbb R^n$、测度为 Lebesgue 测度，使先前球覆盖论证给出的 $A_n$ 可用；$C$ 可以依赖维数，但与 $f,p$ 无关。
\end{remark}
% END SUPPLEMENT: id=v14-ch06-maximal-lp-proof
% SOURCE p=221; printed=206; object=ch06-weak-kernel-introduction; region=lower
第二个应用是一个与\cref{theorem:ch06-kernel-bound}有关的积分算子定理。
% SOURCE p=221; printed=206; object=ch06-weak-kernel; number=6.36; region=lower
\begin{theorem}\label{theorem:ch06-weak-kernel}
设 $(X,\mathscr M,\mu)$、$(Y,\mathscr N,\nu)$ 为 $\sigma$ 有限测度空间，$1<q<\infty$。设 $K$ 为 $X\times Y$ 上可测函数，且存在 $C>0$，使几乎所有 $x\in X$ 满足 $[K(x,\cdot)]_q\leqslant C$，几乎所有 $y\in Y$ 满足 $[K(\cdot,y)]_q\leqslant C$。原书陈述：若 $1\leqslant p<\infty$、$f\in L^p(\nu)$，则积分
\begin{align*}
Tf(x)=\int_Y K(x,y)f(y)\,\mathrm d\nu(y)
\end{align*}
对几乎所有 $x\in X$ 绝对收敛；这样定义的 $T$ 为弱 $(1,q)$ 型，并且对所有满足 $1<p<r<\infty$、$p^{-1}+q^{-1}=r^{-1}+1$ 的 $p,r$，为强 $(p,r)$ 型。更准确地，有与 $K$ 无关的常数 $B_p$，使
\begin{align*}
[Tf]_q&\leqslant B_1C\|f\|_1,\\
\|Tf\|_r&\leqslant B_pC\|f\|_p,\qquad p>1,\quad r^{-1}=p^{-1}+q^{-1}-1>0.
\end{align*}
\end{theorem}
% SUPPLEMENT: id=v14-ch06-weak-kernel-range; parent=ch06-weak-kernel; kind=source-correction+counterexample+hypotheses; origin=independent-derivation
\begin{remark}
\textbf{原文绝对收敛范围的澄清与反例（重构新增）。}原页开头写 $1\leqslant p<\infty$，但其证明下一页实际使用 $1\leqslant p<q'$，其中 $q'=q/(q-1)$。能够由这些条件保证绝对收敛的范围是后者；所列弱型、强型结论的指数均在此范围内，保持不变。不能把开头较宽的字面范围独立当作正确结论。

具体地，取 $X$ 为测度一的单点，$Y=(1,\infty)$ 配 Lebesgue 测度，$K(*,y)=y^{-1/q}$。每列作为单点上的函数，弱 $L^q$ 准范数为 $y^{-1/q}\leqslant1$；每行的水平集长度为 $(a^{-q}-1)_+$，故弱 $L^q$ 准范数也为一。当 $p=q'$ 时，选 $1/p<a\leqslant1$，令
\begin{align*}
f(y)=y^{-1/p}(\log y)^{-a}\chi_{(e,\infty)}(y).
\end{align*}
则 $\int|f|^p=\int_e^\infty y^{-1}(\log y)^{-ap}\,\mathrm dy<\infty$，但
$\int Kf=\int_e^\infty y^{-1}(\log y)^{-a}\,\mathrm dy=\infty$。当 $p>q'$ 时，取 $1/p<\gamma\leqslant1/q'$，令 $f(y)=y^{-\gamma}$；这时 $\gamma p>1$，故 $f\in L^p$，而 $1/q+\gamma\leqslant1$，故 $\int Kf=\infty$。两例均满足原核的全部行列假设，排除了较宽范围。下面保留原证明，并按其实际使用的 $p<q'$ 理解绝对收敛结论。
\end{remark}
% END SUPPLEMENT: id=v14-ch06-weak-kernel-range
% SOURCE p=221,222,223; printed=206-208; object=ch06-weak-kernel-proof; region=bottom-top
\begin{proof}
设 $p',q'$ 分别为 $p,q$ 的共轭指数，则
\begin{align*}
r^{-1}=p^{-1}+q^{-1}-1=p^{-1}-(q')^{-1}=q^{-1}-(p')^{-1},
\end{align*}
所以 $p<q'$ 且 $q<p'$。设 $0\ne f\in L^p$，$1\leqslant p<q'$；给 $f,K$ 乘以常数后，可假设 $\|f\|_p=C=1$。给定稍后选择的 $A>0$，定义
\begin{align*}
E=\{(x,y):|K(x,y)|>A\},\qquad
K_1=(\operatorname{sgn}K)(|K|-A)\chi_E,\qquad K_2=K-K_1,
\end{align*}
并记对应的算子为 $T_1,T_2$。由层积分与截断分布公式，因为 $q>1$，
\begin{align*}
\int_Y|K_1(x,y)|\,\mathrm d\nu(y)
&=\int_0^\infty\lambda_{K(x,\cdot)}(\alpha+A)\,\mathrm d\alpha
\leqslant\int_A^\infty\alpha^{-q}\,\mathrm d\alpha
=\frac{A^{1-q}}{q-1},\\
\int_X|K_1(x,y)|\,\mathrm d\mu(x)&\leqslant\frac{A^{1-q}}{q-1}.
\end{align*}
由\cref{theorem:ch06-kernel-bound}，定义 $T_1f$ 的积分几乎处处收敛，而且
\begin{align}
\|T_1f\|_p\leqslant\frac{A^{1-q}}{q-1}\|f\|_p
=\frac{A^{1-q}}{q-1}.
\label{align:ch06-weak-kernel-large}
\end{align}
同样，因为 $q<p'$，
\begin{align*}
\int_Y|K_2(x,y)|^{p'}\,\mathrm d\nu(y)
&=p'\int_0^A\alpha^{p'-1}\lambda_{K(x,\cdot)}(\alpha)\,\mathrm d\alpha\\
&\leqslant p'\int_0^A\alpha^{p'-1-q}\,\mathrm d\alpha
=\frac{p'A^{p'-q}}{p'-q}.
\end{align*}
因此由 Hölder 不等式，定义 $T_2f(x)$ 的积分收敛（原文写为对每个 $x$；关于例外行及 $p'=\infty$ 的解释见后注），且
\begin{align}
\|T_2f\|_\infty\leqslant
\left[\frac{p'A^{p'-q}}{p'-q}\right]^{1/p'}\|f\|_p
=\left(\frac rq\right)^{1/p'}A^{q/r}.
\label{align:ch06-weak-kernel-small}
\end{align}
于是 $Tf=T_1f+T_2f$ 几乎处处定义良好。

给定 $\alpha>0$，由分布函数的和估计，
\begin{align*}
\lambda_{Tf}(\alpha)\leqslant
\lambda_{T_1f}(\alpha/2)+\lambda_{T_2f}(\alpha/2).
\end{align*}
按\cref{align:ch06-weak-kernel-small}取
\begin{align*}
A=(\alpha/2)^{r/q}(q/r)^{r/(qp')},
\end{align*}
便有 $\|T_2f\|_\infty\leqslant\alpha/2$，因而第二项为零。由\cref{align:ch06-weak-kernel-large}及 Chebyshev 不等式，
\begin{align*}
\lambda_{Tf}(\alpha)
&\leqslant\lambda_{T_1f}(\alpha/2)
\leqslant\left(\frac{2\|T_1f\|_p}\alpha\right)^p
\leqslant\left[\frac{2A^{1-q}}{(q-1)\alpha}\right]^p\\
&=\frac{2^{p-(1-q)pr/q}}{(q-1)^p}
\left(\frac qr\right)^{(1-q)pr/(qp')}
\alpha^{-p+(1-q)pr/q}
=C_p\left(\frac{\|f\|_p}\alpha\right)^r.
\end{align*}
最后一步用了 $\|f\|_p=1$ 与
\begin{align*}
\frac{(1-q)pr}q-p
=p\left(-\frac r{q'}-1\right)=-p\frac rp=-r.
\end{align*}
简单的齐次性论证恢复任意 $\|f\|_p$，故 $T$ 为弱 $(p,r)$ 型；特别在 $p=1$ 时，为弱 $(1,q)$ 型。

最后，给定 $p\in(1,q')$，取 $\widetilde p\in(p,q')$，并定义 $\widetilde r^{-1}=\widetilde p^{-1}-(q')^{-1}$。$T$ 为弱 $(1,q)$ 型与弱 $(\widetilde p,\widetilde r)$ 型，所以由 Marcinkiewicz 定理，$T$ 为强 $(p,r)$ 型。
\end{proof}
% SUPPLEMENT: id=v14-ch06-weak-kernel-completion; parent=ch06-weak-kernel; kind=omitted-case+proof-detail+motivation+hypotheses; origin=independent-derivation
\begin{remark}
\textbf{绝对收敛、端点与条件说明（重构新增）。}先固定一个 $A>0$。对 $K_1$ 应用的核定理不仅给出 $T_1f$ 的范数界，也给出 $\int|K_1(x,y)f(y)|\,\mathrm d\nu(y)<\infty$ 几乎处处。对 $K_2$，上述 Hölder 估计给出同样的绝对收敛，因此原积分绝对收敛，不依赖后来选用哪个阈值。这也保证线性及不同表示的相容性；联合可测性与 $\sigma$ 有限性通过非负积分和正负部分保证输出的可测性。

行界只对几乎所有 $x$ 假定成立，所以 $T_2$ 的收敛结论首先也只对这些行成立；原文的“每个 $x$”须限于这个满测集。若希望处处表述，可以先在行界失败的可测零集上把 $K$ 改为零；这不改变算子的几乎处处类或列积分估计。没有必要对未满足行界的任意代表强加逐点收敛。

当 $p=1$ 时，$p'=\infty$，原证明中的 $p'$ 次方积分不能照字面使用。此时直接由 $|K_2|\leqslant A$ 得
\begin{align*}
\int|K_2(x,y)f(y)|\,\mathrm d\nu(y)\leqslant A\|f\|_1=A.
\end{align*}
取 $A=\alpha/2$，用 $T_1$ 的 $L^1$ 界，
\begin{align*}
\lambda_{Tf}(\alpha)\leqslant\frac2\alpha\|T_1f\|_1
\leqslant\frac{2^q}{q-1}\alpha^{-q}.
\end{align*}
所以 $B_1=2(q-1)^{-1/q}$ 可用。$1<p<q'$ 时 $p'<\infty$，原计算全部有效；恒等式 $p'/(p'-q)=r/q$ 及 $1-q/p'=q/r$ 直接由指数关系得到。原证的 $C_p$ 只依赖 $p,q$，恢复缩放后弱型界为 $[Tf]_r\leqslant C_p^{1/r}C\|f\|_p$，因为先对 $K/C$、$f/\|f\|_p$ 作归一化；若输入为零，绝对积分也为零，无须除以其范数。

最后的插值也须核实指数。对给定 $1<p<q'$，可取 $\widetilde p=(p+q')/2$，从而 $1< p<\widetilde p<q'$；令 $t$ 由 $1/p=1-t+t/\widetilde p$ 确定，则 $0<t<1$，且
\begin{align*}
\frac{1-t}q+\frac t{\widetilde r}
=\frac{1-t}q+t\left(\frac1{\widetilde p}-\frac1{q'}\right)
=\frac1p-\frac1{q'}=\frac1r.
\end{align*}
端点满足 $1\leqslant q$、$\widetilde p<\widetilde r$、$q<\widetilde r$，所以 Marcinkiewicz 的全部指数条件成立。输入和空间上，每一项积分都绝对收敛，故 $T$ 在和空间上确为线性，从而次线性。得到的强型常数仍只依赖 $p,q$（此处已固定 $\widetilde p$ 的选择），不依赖 $K$。

动机是把只有弱行列界的核切成具有强 $L^1$ 行列界的大值部分，以及行 $L^{p'}$ 可控的小值部分；前者交给核有界性，后者交给 Hölder。$q>1$ 保证大值尾积分 $\int_A^\infty a^{-q}\,\mathrm da$ 收敛；$q<\infty$ 使 $q'$ 为严格大于一的有限上端点。$p<q'$ 等价于 $p'>q$，保证小值积分在零端收敛，也是 $r>0$ 的要求。行、列两种弱界一起控制 $K_1$；统一常数使最终范数界不依赖行列。$\sigma$ 有限性沿用前面核定理与乘积积分的条件，不是由单个截面不等式自动产生的。
\end{remark}
% END SUPPLEMENT: id=v14-ch06-weak-kernel-completion
\end{document}
